\documentclass{article}
\usepackage{graphicx} % Required for inserting images
\usepackage{tikz}
\usepackage{physics}
\usepackage{gensymb}
\usepackage{amsmath, amsfonts, mathtools, amsthm, amssymb, mathrsfs}
\usepackage{titlesec}
\usepackage{fancyhdr}
\usepackage{hyperref}
\usepackage{float}
\usepackage{booktabs}
\usepackage{enumitem}
\usepackage{subcaption}
\usepackage{multicol}
\usepackage{import}
\usepackage[usenames,dvipsnames]{xcolor}
\usepackage[a4paper, margin=1in]{geometry}
% The configuration for boxes and parts of the overall idea come from Pingbang Hu: https://github.com/sleepymalc
% Use a relative path (forward slashes or "..") to avoid backslash-escapes
\subimport{../}{header.tex}

\title{\textbf{MATH 281 Applied Boundary Value Problems}}
\author{Autumn Weaver}
\date{Fall 2026}

\begin{document}

\maketitle

\tableofcontents{\textbf{}}

\pagebreak

\pagestyle{fancy}
\fancyhead{}
\fancyhead[L]{\thetitle}
\fancyhead[R]{Autumn Weaver}

\section{Heat Equation}
\begin{definition}
A \textbf{partial differential equation} (PDE) is an equation 
containing partial derivatives. 
\end{definition}

\subsection{Conduction of Heat in a 1-D Rod}
\begin{definition}
The \textbf{thermal energy density} \(e(x,t)\) is the amount of thermal energy per 
unit volume. 
\end{definition}

If we have 1 dimensional rod with cross sectional area \(A\) and length \(L\) and a slice on the rod 
from \(x\) to \(x+\Delta x\), then the heat energy of this slice is:
\[
    \text{heat energy}=e(x,t)A\Delta x
\]  

If there is no lateral heat flow on the rod, then the rate of change of the heat energy of 
this slice must be the sum of the next heat flux into the slice and the heat generate in the slice:
\[
    \frac{\partial }{\partial t}[e(x,t)A\Delta x]\approx \Phi(x, t)A-\Phi(x+\Delta x, t)A+Q(x, t)A\Delta X
\]
Where \(Q(x, t)\) is the heat energy generated per unit volume.


We can easily simplify this by diving by \(A \Delta x\) and taking the \(\lim_{\Delta x \rightarrow 0}\)  to get the following PDE:
\[
    \frac{\partial e}{\partial t}=-\frac{\partial \Phi}{\partial x}+Q
\]

We can also derive this equation by considering a finite slice from \(x=a\) to \(x=b\). The rate of change of heat energy of this slice is: 
\[
    \frac{d}{dt}\int_a^b e dx=\Phi(a,t)-\Phi (b,t)+\int_a^b Q dx
\]

This is also known as the \textbf{integral conservation law}. 

We can convert this to a single integral by noting that \(\frac{d}{dt}\int_a^b e(x, t)=\int_a^b\frac{\partial e(x, t)}{t}\) and that 
\(\Phi(a, t)-\Phi(b, t)=-\int_a^b \frac{\partial \Phi}{\partial x}dx\), resulting in: 
\[
    \int_a^b (\frac{\partial e}{\partial t}+\frac{\partial \Phi}{\partial x}-Q)dx=0
\]

Since this must be true for any arbitrary \(a\) and \(b\), the integrand must be zero:
\[
    \frac{\partial e}{\partial t}=-\frac{\partial \Phi}{x}+Q
\]

Because we also describe the temperatures of materials instead of heat energy, it is useful to convert our equations. We let:
\begin{align*}
    u(x, t)&\equiv\text{temperature} \\
    c&\equiv\text{specific heat capacity} \\
    \rho&\equiv\text{density}
\end{align*}



\end{document}